\section{图预处理与精确计数流程}
\subsection{图清洗与编号映射}
首先收集输入端点，将非自环边规范为 $(\min(u,v),\max(u,v))$，排序后合并重边。原始顶点编号排序去重，得到连续下标。编号映射保留原编号之间的大小关系，后续数组使用内部下标。根据清洗后的无向边计算每个顶点的度数 $d(u)$。

\subsection{度数优先定向与唯一计数}
定义全序 $u\prec v$：当 $d(u)<d(v)$，或 $d(u)=d(v)$ 且 $u<v$ 时成立。每条无向边从全序中靠前的端点指向靠后的端点，得到定向边集 $E^+$。令 $F(u)=\{v:(u,v)\in E^+\}$，则
\begin{equation}
T=\sum_{(u,v)\in E^+}|F(u)\cap F(v)|.
\label{eq:triangle}
\end{equation}
对于任意三角形，按全序记其三个顶点为 $a\prec b\prec c$。顶点 $c$ 同时属于 $F(a)$ 与 $F(b)$，因此该三角形在边 $(a,b)$ 的交集中被计入。另两条边不再为同一三元组贡献计数，故每个三角形恰计一次，无需除以3或6。边定向只是计算内部安排，不改变无向计数语义。

\subsection{向前CSR与二分交集}
向前邻居使用CSR存储，各行内部按顶点编号排序。需要区分：边方向由 $(d(u),u)$ 决定，行内排序依据编号。CSR使用约 $8(|V|+1)+4|E|$ 字节，另计任务元数据及部分和；此处每条清洗后的无向边只存一个方向。

对于一条边 $(u,v)$，遍历较短向前表中的元素，在较长表中二分查找。若两表长度为 $a,b$，搜索工作量可用 $\min(a,b)\log_2(\max(a,b)+1)$ 粗略估计，但实际耗时还受提前命中、访存与线程执行影响。空表可直接返回0。CPU参考使用双指针交集，辅助核验GPU二分实现。

\subsection{基础GPU任务划分与归约}
每条定向边由一个GPU线程完整处理；线程按全局下标跨步领取后续边。额外源点数组记录每条CSR边的源点，终点直接使用邻居数组。线程先累加局部64位计数，通过CUB完成块内归约与最终设备级归约，再将总数回传主机。实验大图使用864个线程块、每块256个线程，这是逻辑启动规模，不表示所有线程同时驻留。
